高被引用論文や引用指標の研究は、近年も幅広い観点から進展している。Liら\cite{Li1}は2300万論文を対象に長期にわたる引用の特徴を分析しグループ構造を明らかにした。また、Yuら\cite{Yu1}は社会心理学分野のトップ被引用論文100本を対象とした構造解析により、主要な出版年・国際協力・ジャーナルの影響構造を明らかにし、被引用論文の研究ネットワーク構造に関する知見を提供している。 雑誌レベルに注目したものとして、Leydesdorff\cite{Ley1,Ley2,Ley3,Ley4}は雑誌間の関係マップを構築・調査した。天野ら\cite{Ama1,Ama2}は医学生物学雑誌の大規模クラスタリングを行い、分野間の引用関係を明らかにした。こうした研究は、高被引用論文あるいはハイインパクト誌がどのような条件下で生まれ、どのような影響経路を持つかをマクロ的な視点で捉える試みである。一方で、論文レベルの被引用に寄与する構造的因子の体系化も進んでおり、共同研究形態・出版戦略・研究資源配分などの多要因が被引用に影響する可能性が示唆されている。Hillら\cite{Hil1}は、研究者の研究テーマの転換戦略はむしろその後の影響を減少させる可能性を指摘した。Linら\cite{Lin1}は、被引用論文の全文データを用いて、引用の位置・文脈・引用者との類似性が経時的に変化する現象を観察し、単なる引用数では捉えきれない引用の役割の変容を指摘している。Koushaら\cite{Ko1}は、情報科学分野の論文に対し、抄録の「読みやすさ」と論文の「質」が、被引用数にどのように影響するかを調査した。以上紹介した研究は主として被引用性にかかるさまざまな属性の特徴を明らかにするものであるが、高被引用論文に共通する長期的構造特性そのものを多層的に捉える研究は限定的である。本研究は、引用寿命・雑誌構造・テーマ構造などの要素を統合的に分析することで既存研究とは異なる視点を提供する。
